\section{Arbitrary-phase nondegeneracy and a local Hessian obstruction}\label{sec:nondegeneracy}
We study the second-order cohomological variance \(\Gamma_A(g)=\sigma_A^2(\{g,g\circ A^{\mathsf T}\}/(4\pi))\) for hyperbolic two-torus automorphisms. For the fixed Pell matrix we give a finite, independently replayable proof that it is positive for every nonconstant real trigonometric polynomial of frequency radius at most eight. All 288 real parameters are arbitrary: evenness and Fourier-phase alignment are not assumed. An exact finite quantum variance identity connects this certificate to a period-averaged observable, with an explicit non-aliasing hypothesis. A separate analytic theorem detects nonzero variance from the two hyperbolic curvatures at a fixed critical point. It implies nondegeneracy, without a bandwidth restriction, on the cone of even symbols with one-sign cosine coefficients, for every hyperbolic integral determinant-one matrix.

These are project proofs and exact finite certificates, not a claim of world priority or positive BGS universality. The Fourier-orbit formula for asymptotic variance is established theory; we prove the specialized identities needed here. The certificate is stronger than sampling real coefficients but is not a Lean formalization. The unbounded-radius mixed-sign problem remains unresolved by these arguments. No uniform long-time quantum comparison is inferred from a positive classical variance alone.

\subsection{Normalization and the cohomological question}\label{nondegeneracy:normalization-and-the-cohomological-question}

Let \(A\in\mathrm{SL}(2,\mathbb Z)\) be hyperbolic, \(S=A^{\mathsf T}\), and \(\mathcal C_A f=f\circ S\). On the unit torus put \[
 e_v(x)=e^{2\pi i v\cdot x},\qquad b_v=e_v+e_{-v},\qquad
 \{f,h\}=f_{x_1}h_{x_2}-f_{x_2}h_{x_1}.
\] For a real smooth mean-zero function \(h\), define \[
 \sigma_A^2(h)=\sum_{\mathcal O}
       \left|\sum_{v\in\mathcal O}\widehat h(v)\right|^2,
 \qquad
 \Gamma_A(g)=\sigma_A^2\!\left(\frac{\{g,\mathcal C_Ag\}}{4\pi}\right),
 \tag{ND.1}
\] where \(\mathcal O\) ranges over \textbf{oriented} \(A\)-orbits in \(\mathbb Z^2\setminus\{0\}\). An orbit and its negative are distinct: a nonzero integer vector cannot be periodic under \(-A^j\) for a hyperbolic \(A\). For real \(h\) they contribute conjugate sums. For even \(h=\sum a_v b_v\), (ND.1) therefore reduces to twice the sum of the squared cosine sums on signed orbits. Constants in \(g\) affect neither the bracket nor \(\Gamma_A\).

For finite Fourier support, character orthogonality gives exactly \[
 \sum_{n\in\mathbb Z}\int h(x)h(S^n x)\,dx
   =\sum_{\mathcal O}\left|\sum_{v\in\mathcal O}\widehat h(v)\right|^2.
 \tag{ND.2}
\] The double absolute sum is bounded by \(\|\widehat h\|_1^2\), so the same identity holds for smooth \(h\) by absolute convergence. This is the usual Green--Kubo/orbit-sum variance, not the instantaneous \(L^2\) norm. Such spectral and orbit-sum descriptions are prior theory, including Cohen--Conze {[}1{]}.

For a finite polynomial \(h\), \(\sigma_A^2(h)=0\) if and only if \(h=u-\mathcal C_Au\) for a finite polynomial \(u\). Indeed, on each orbit, let \(h_j\) be the finite coefficient sequence. If its sum is zero, the partial sums \(u_j=\sum_{k\leq j}h_k\) have finite support and \(h_j=u_j-u_{j-1}\). Conversely a finite difference has zero sum. Consequently proving only \(\{g,\mathcal C_Ag\}\ne0\) does not settle the nondegeneracy question.

For Sections 2--3 fix \[
 T=\begin{pmatrix}3&2\\4&3\end{pmatrix},\qquad
 \lambda=3+2\sqrt2.
\] Let \(\mathcal P_R\) be the real mean-zero trigonometric polynomials supported in \(\|v\|_\infty\leq R\). Its real dimension is \(4R(R+1)\).

\subsection{An exact finite theorem for arbitrary Fourier phases}\label{nondegeneracy:an-exact-finite-theorem-for-arbitrary-fourier-phases}

\textbf{Theorem 1 (the full radius-eight box).} For every nonzero \(g\in\mathcal P_8\), \[
                         \boxed{\Gamma_T(g)>0.}             \tag{ND.3}
\] In particular (ND.3) holds at every smaller integer radius. There is no restriction on the signs or phases of the Fourier coefficients. This proves that the second-order bracket is not a finite Fourier coboundary throughout this 288-dimensional real space. The assertion is not extrapolated to all radii.

\subsubsection{2.1 Exact polynomial equations}\label{nondegeneracy:exact-polynomial-equations}

Choose \[
 V_R=\{(a,b):0\leq a\leq R,\ -R\leq b\leq R,
                         \ a>0\text{ or }b>0\}.
\] Write, uniquely, \[
 g=\sum_{v\in V_R}(z_v e_v+\overline{z_v}e_{-v}),\qquad
 z_v\in\mathbb C,\qquad \|g\|_2^2=2\sum_v|z_v|^2.          \tag{ND.4}
\] There are \(n_R=2R(R+1)\) complex variables. The exact Fourier bracket is \[
 \widehat{\{g,\mathcal C_Tg\}/(4\pi)}(k)
       =-\pi\sum_{u+Tv=k}\det(u,Tv)\widehat g(u)\widehat g(v).
 \tag{ND.5}
\] Combining the coefficients along one of the two orientations of each signed orbit gives integer-coefficient quadratic polynomials \(Q_j(z,\bar z)\) with \[
                   \Gamma_T(g)=2\pi^2\sum_j|Q_j(z,\bar z)|^2. \tag{ND.6}
\] Thus (ND.3) is the assertion that all the equations \(Q_j=0\) have only the zero solution subject to the genuine complex-conjugation relation. The latter is not replaced by treating \(z\) and \(\bar z\) as independent complex variables.

\subsubsection{2.2 Why the orbit classification terminates and preserves phases}\label{nondegeneracy:why-the-orbit-classification-terminates-and-preserves-phases}

For \(v=(a,b)\) put \(\beta=b+\sqrt2 a\), \(\beta'=b-\sqrt2 a\). Then \[
 E(T^jv):=2(T^jv)_1^2+(T^jv)_2^2
       =\tfrac12(\beta^2\lambda^{2j}+\beta'^2\lambda^{-2j}). \tag{ND.7}
\] Both coefficients are positive for \(v\ne0\). This is strictly convex as a function of real \(j\), tends to infinity at both ends, and has at most two adjacent integer minimizers. Integer neighbour descent reaches the global minimum. At equal energies we use a fixed lexicographic rule, including the choice of a sign for the vector. Energy decreases in the non-tie steps, so the algorithm cannot loop indefinitely.

Crucially, the program records the orientation bit \(\epsilon\) in \(w=\epsilon T^jv\). A negative orientation contributes to the conjugate orbit, not a second independent copy of the same complex coefficient. Adjacent minimum ties, for example the orbit containing \((1,-2)\), are retained as explicit normalization tests. Omitting this bit can produce an incorrect arbitrary-phase variance even when all even-symbol tests pass.

The independent C++ constructor works in the quadratic ring with vector \(b+a\sqrt2\) and enumerates an orbit window \(-16\leq j\leq16\). It fails if a minimizer is at the boundary. Interior minimizers have the global meaning just established by (ND.7); the finite window is not an unproved coverage assumption. All 1,475 nonzero radius-eight equations agree exactly, term by term, with the independently constructed Python equations.

\subsubsection{2.3 A small proof language and its soundness}\label{nondegeneracy:a-small-proof-language-and-its-soundness}

The proof records known zero coefficients, known zero products, and temporary nonzero assumptions. Only the following four rules are admitted.

\textbf{P (product).} After terms killed by established facts are deleted, an equation has one cross monomial \(a z_i^{\epsilon}z_j^{\eta}=0\), with \(a\ne0\) and \(i\ne j\). Record \(z_i z_j=0\), equivalently that at least one factor is zero. Here the exponent sign denotes the variable or its conjugate, not an inverse.

\textbf{S (squares).} A sole nonzero multiple of \(z_i^2\), \(\bar z_i^2\), or \(|z_i|^2\) implies \(z_i=0\). A sum \(\sum_i a_i|z_i|^2=0\) with all nonzero \(a_i\) of one sign implies that every listed variable vanishes. A sum of ordinary complex squares is \emph{not} admitted by the second rule.

\textbf{D (nonzero division).} A recorded zero product and a currently assumed nonzero factor imply that the other factor is zero. Division without this explicit assumption is rejected.

\textbf{F (refutation).} Temporarily assume \(z_i\ne0\), replay the supplied child proof with inherited established facts, and require an actual contradiction. If the child succeeds, conclude \(z_i=0\) outside the branch. An empty child or a branch that merely makes partial progress is rejected.

Each rule is an elementary implication over \(\mathbb C\) with its usual conjugation. Induction on the sequential deductions and refutations proves soundness of a complete certificate. The search need not be complete: failure to eliminate all variables is reported as UNKNOWN, never as a counterexample. The validator replays the proof and does not trust the generator's status.

The radius-eight certificate contains 5,704 elementary deductions: 3,293 P rules, 1,077 S rules, 1,229 D rules and 105 F rules. Every one of its 144 complex variables is eliminated. Refutations in this particular certificate have depth one. Both Python and C++ independently reconstruct the coefficient system and replay the deductions. Four malformed Python objects and two malformed C++ instruction streams are required to fail closed. This completes a finite computer-assisted proof of (ND.3), with the analytic normalization (ND.5)--(ND.7) and logical soundness above as explicit dependencies.

\begin{longtable}[]{@{}rrrrrrr@{}}
\toprule\noalign{}
Radius & Real dimension & Equations & P & S & D & F \\
\midrule\noalign{}
\endhead
\bottomrule\noalign{}
\endlastfoot
1 & 8 & 23 & 0 & 3 & 0 & 0 \\
3 & 48 & 212 & 81 & 35 & 16 & 4 \\
8 & 288 & 1,475 & 3,293 & 1,077 & 1,229 & 105 \\
\end{longtable}

The radius-eight coefficient table has SHA-256 \path{7462bc1a522a57741028efd6642710fa50aa271b27f21bfc6a683f578aee53f4}. The plain-text proof has SHA-256 \path{ed26c33904551e11a5185d55702d781b4f2555d7a14153ceaf6e8a0b5f71d94a}. The delivered manifest provides all source and result hashes. Hash agreement identifies a proof object; it does not by itself prove that object sound.

The certificate objects are in the repository. The plain-text proof itself
(the 5,704 P/S/D/F deductions, with rule counts 3,293/1,077/1,229/105) is
\path{evidence/executions/20261006/corpus-revision/final/checks/certificates-mpuvq_28/resonance_nondegeneracy_radius8.proof}
(git blob \path{17a70f2b42166c0e8a339050d38f3b1883483720}; its SHA-256 is
the plain-text proof hash above). The accompanying certificate record,
stating radius $8$, real dimension $288$, $144$ complex variables and
$1{,}475$ orbit equations with the coefficient-table hash above, is
\path{evidence/executions/20261006/corpus-revision/final/checks/certificates-mpuvq_28/resonance_nondegeneracy_radius8.json}
(git blob \path{dcbe0a7d1a208c6596d37fd4178d20f9cbd07a52}). The replay
receipt
\path{evidence/executions/20261006/corpus-revision/final/checks/certificates-mpuvq_28/resonance_certificate_replay.json}
(git blob \path{3621e53121a776596d936829b7a6b75bd97975d2}) records PASS
with the same rule counts and both hashes. The generator and the replay
checker are \path{computations/verification/resonance_nondegeneracy_certificate.py}
(git blob \path{6503d86b52b713bfb042031a7e5a97e6f62e8cc8}) and
\path{computations/verification/check_resonance_certificate.py}
(git blob \path{67b715a3888baa2da7d288a5e3d577ebc3d501c1}).

\subsection{A human-readable arbitrary-phase example and a quantitative margin}\label{nondegeneracy:a-human-readable-arbitrary-phase-example-and-a-quantitative-margin}

For \(R=1\), in the order \((0,1),(1,-1),(1,0),(1,1)\), three of the exact chosen orbit coefficients are \[
 Q_{(-4,-4)}=-4\bar z_2^2,\quad
 Q_{(-2,-4)}=2\bar z_0^2+4|z_2|^2,\quad
 Q_{(0,-2)}=2|z_1|^2+2|z_3|^2.                              \tag{ND.8}
\] Their successive vanishing immediately forces all four variables to vanish. This is not restricted to a real coefficient vector.

A useful, though not sharp, bound follows without the larger certificate: \[
              \Gamma_T(g)\geq\frac{8\pi^2}{17}\|g\|_2^4,
                         \qquad g\in\mathcal P_1.          \tag{ND.9}
\] Indeed, put \(x=|z_0|^2\), \(y=|z_2|^2\), \(w=|z_1|^2+|z_3|^2\). Equation (ND.8) and the reverse triangle inequality give \(\Gamma_T(g)\geq8\pi^2[(x-2y)^2+4y^2+w^2]\). Cauchy--Schwarz gives \((x+y+w)^2\leq(17/4)[(x-2y)^2+4y^2+w^2]\), and (ND.4) proves (ND.9). The sharper even-symbol constant from the earlier project is not silently transferred to arbitrary Fourier phases.

\subsubsection{A constructive margin for the radius-eight certificate}\label{nondegeneracy:a-constructive-margin-for-the-radius-eight-certificate}

Compactness already gives some \(c_8>0\) with \(\Gamma_T(g)\geq c_8\|g\|_2^4\). There is also a completely explicit, deliberately very poor lower bound: \[
 \boxed{\displaystyle
 \Gamma_T(g)\geq \frac{\pi^2}{41472}\,2^{-2^{715}}\|g\|_2^4,
                         \qquad g\in\mathcal P_8.}          \tag{ND.10}
\] This number has no useful interpretation as a physical crossover estimate; its purpose is to remove an unspecified compactness constant. It is not claimed to approximate the sharp minimum.

Here is the quantitative compilation rule. Normalize \(\max_i|z_i|=1\) and suppose \(|Q_j|\leq\delta\). Replace an exact zero fact by an upper bound \(2^L\delta^{p-kb}\) for its absolute value. The symbol \(b\) appears only under a temporary assumption \(|z_i|\geq\delta^b\). Every deleted monomial is bounded by a known coefficient bound times \(|z_j|\leq1\), or by a known product bound. In the residual row, use the triangle inequality, the minimum of the base exponents, the maximum of their \(b\) slopes, and a power-of-two upper bound on the sum of absolute integer coefficients. P records the resulting product bound. S takes a square root, rounding its power-of-two coefficient upward. D subtracts \(b\) from the exponent using the explicit nonzero lower bound.

At the end of a refutation child, the pivot has bound \(2^L\delta^{p-kb}\) with \(p>0\). Choose \(b=p/[2(k+1)]\) and require \(\delta\leq2^{-B}\) with \(B\geq2(L+1)/p\). Then this bound is strictly less than \(\delta^b\), contradicting the temporary assumption. Therefore the pivot has the unconditional bound \(|z_i|<\delta^b\). Proceed through the root proof and finally enlarge \(B\) so that every coefficient is at most \(1/2\). This contradicts the normalization. All exponents and comparisons in the compiler are exact rational arithmetic; there is no numerical interval test on the coefficients of \(g\).

For the maintained radius-eight proof the generated integer \(B\) has 714 binary digits, so the simple upper bound \(B\leq2^{714}\) suffices. The preceding argument gives \(\max_j|Q_j|>2^{-B}\) on \(\max_i|z_i|=1\). Since \(\|g\|_2^4\leq4n_R^2(\max_i|z_i|)^4\), (ND.6) yields \(\Gamma_T(g)\geq[\pi^2/(2n_R^2)]2^{-2B}\|g\|_2^4\); with \(n_8=144\) this implies (ND.10). The exact, much less conveniently printed integer and every intermediate rational bound are retained in the margin report. The original margin report was not deposited; it has been reconstructed from the deposited certificate by an exact-rational-arithmetic implementation of the compilation rule above (\path{computations/verification/resonance_nondegeneracy_margin.py}), and the reconstructed report is deposited at \path{evidence/executions/20261006/corpus-revision/final/checks/certificates-mpuvq_28/resonance_nondegeneracy_margin.json}. The reconstructed integer has 712 binary digits, so \(B\leq2^{714}\) is confirmed and the explicit constant in (ND.10) is recheckable from deposited materials; two independent implementations of the compiler produce the identical reconstructed integer. The unconditional compactness bound \(\Gamma_T(g)\geq c_8\|g\|_2^4\) for some \(c_8>0\) does not depend on the margin report.

\subsection{An invariant Hessian obstruction for every hyperbolic matrix}\label{nondegeneracy:an-invariant-hessian-obstruction-for-every-hyperbolic-matrix}

The finite certificate settles a large but fixed frequency box. A separate analytic criterion has no bandwidth restriction and is not specific to T. Write \[
 A=\begin{pmatrix}a&b\\c&d\end{pmatrix},\quad t=a+d,\quad
 \Delta=t^2-4>0,\quad
 B_A=\begin{pmatrix}c&(d-a)/2\\(d-a)/2&-b\end{pmatrix}.
\] The integral quadratic form \[
 q_A(v)=v^{\mathsf T}B_Av=\det(v,Av)
\] is A-invariant. Define, for a symmetric matrix H, \[
 \mathcal E_A(H)=(\operatorname{tr}(B_AH))^2+\Delta\det H,
 \qquad
 \mathcal L_A=c\partial_{x_1}^2+(d-a)\partial_{x_1}\partial_{x_2}
                                  -b\partial_{x_2}^2.       \tag{ND.11}
\]

\textbf{Theorem 2 (fixed critical-point obstruction).} Let \(g\) be real smooth. If \(x_0\) is a fixed point of \(S=A^{\mathsf T}\) and \(dg(x_0)=0\), then \[
 \boxed{\Gamma_A(g)=0\quad\Longrightarrow\quad
                   \mathcal E_A(\operatorname{Hess}g(x_0))=0.} \tag{ND.12}
\] In particular, if the two pure second derivatives in the stable and unstable coordinate directions are both nonzero at that point, \(\Gamma_A(g)>0\). Every even \(g\) is critical at the fixed point zero, but evenness is not a hypothesis of Theorem 2. A nondegenerate extremum at a fixed point is a sufficient condition, not a necessary one.

\textbf{Proof.} Put \(h=\{g,\mathcal C_Ag\}/(4\pi)\). At a fixed point, \(e_{Av}(x_0)=e_v(x_0)\), and \(q_A(Av)=q_A(v)\). Grouping the absolutely convergent twice-differentiated Fourier series by oriented orbits gives \[
 \mathcal L_Ah(x_0)
 =-4\pi^2\sum_{\mathcal O}q_A(v_{\mathcal O})e_{v_{\mathcal O}}(x_0)
                       \sum_{v\in\mathcal O}\widehat h(v).
 \tag{ND.13}
\] If the variance is zero, every orbit sum is zero, so (ND.13) vanishes. No regularity of a putative transfer function has been assumed.

Only the quadratic Taylor jet of g contributes to the second derivative of its bracket at a critical point. For \(g_2=(H_{11}x_1^2+2H_{12}x_1x_2+H_{22}x_2^2)/2\), direct differentiation yields \[
 \mathcal L_A\{g_2,g_2\circ A^{\mathsf T}\}
                         =t\mathcal E_A(H).                \tag{ND.14}
\] Equivalently, diagonalize S in an area-one eigenbasis. If its eigenvalues are \(\mu,\mu^{-1}\), set \(\gamma=\mu-\mu^{-1}\); then \(\gamma^2=\Delta\), \(\mathcal L_A=-\gamma\partial_u\partial_s\), and \[
 \mathcal E_A(H)=\Delta H_{uu}H_{ss}.
 \tag{ND.15}
\] The mixed quadratic coefficient of the bracket is \((\mu^{-2}-\mu^2)H_{uu}H_{ss}\), giving (ND.14). This remains true for either sign of the trace. Since \(t\ne0\), (ND.13)--(ND.14) prove (ND.12). End of proof.

The identity (ND.14) was also checked as a polynomial identity with formal matrix entries and Hessian entries. Before the determinant-one specialization, its discriminant is \(t^2-4\det A\). Two hundred integer instances provide a separate regression, not the general proof.

The converse of (ND.12) is false. For the Pell matrix, \[
 g=4b_{(1,0)}-b_{(2,0)},\qquad \operatorname{Hess}g(0)=0,
 \qquad \Gamma_T(g)=25600\pi^2>0.                            \tag{ND.16}
\] The exact Fourier computation is retained as a negative control. A vanishing Hessian obstruction is therefore not reported as a zero of the variance.

\subsubsection{A quantitative smooth version}\label{nondegeneracy:a-quantitative-smooth-version}

The same obstruction has an explicit lower bound when one controls a weighted Fourier tail. Put \(C_A=|c|+|d-a|+|b|>0\), let \(s>0\), and set \[
 \mathcal H_s(h)=\sum_v|\widehat h(v)|(1+|v|_\infty)^{s+2},
 \qquad
 J=\frac{|\mathcal L_Ah(x_0)|}{4\pi^2}
   =\frac{|t\mathcal E_A(H)|}{16\pi^3}.
\] If \(J>0\), choose an integer \[
 R\geq\max\{1,(2C_A\mathcal H_s(h)/J)^{1/s}\}.
\] Then \[
                  \Gamma_A(g)\geq\frac{J^2}{32C_A^2R^6}>0. \tag{ND.17}
\] To prove it, choose each orbit's representative of minimum frequency radius. There are at most \(8R^2\) oriented orbits meeting the radius-R box, each with \(|q_A|\leq C_AR^2\). Cauchy--Schwarz bounds their part of (ND.13) divided by \(4\pi^2\) by \(\sqrt8 C_AR^3\sqrt\Gamma\). For all other orbits, every frequency has radius greater than R, and the sum is bounded by \(C_AR^{-s}\mathcal H_s(h)\leq J/2\). Rearranging proves (ND.17). This bound is for a fixed regular symbol with a nonzero obstruction. It is not uniform over all mixed-sign Fourier polynomials.

\subsection{A one-sign Fourier theorem without a bandwidth restriction}\label{nondegeneracy:a-one-sign-fourier-theorem-without-a-bandwidth-restriction}

\textbf{Theorem 3.} Let \(A\in\mathrm{SL}(2,\mathbb Z)\) be hyperbolic. If the nonconstant coefficients of a real even smooth mean-zero function \(g=\sum_v c_vb_v\) all have the same sign and at least one is nonzero, then \[
                             \Gamma_A(g)>0.                 \tag{ND.18}
\] This is a statement about cosine coefficients, not pointwise nonnegativity of a mean-zero function. There is no restriction on the number of signed orbits.

\textbf{Proof.} Change the overall sign if necessary, so \(c_v\geq0\), and put \[
 M=\sum_v c_vvv^{\mathsf T}=-\operatorname{Hess}g(0)/(8\pi^2).
\] Polarization of \(q_A\) gives the exact discriminant identity \[
 B_A(u,v)^2-q_A(u)q_A(v)=\frac\Delta4\det(u,v)^2.            \tag{ND.19}
\] Expanding \(\mathcal E_A(M)\) and \(\det M\) shows \[
       \mathcal E_A(M)\geq\left(\sum_v c_v|q_A(v)|\right)^2>0.
 \tag{ND.20}
\] For clarity, each cross coefficient in the difference between the two sides is zero or positive: if \(q_A(u)q_A(v)\geq0\), its extra term is \(\Delta\det(u,v)^2\); if the signs differ, it is \(4B_A(u,v)^2\) by (ND.19). The diagonal terms agree. Also \(q_A(v)\ne0\) for a nonzero integer v: otherwise A would have a rational eigenline and a rational algebraic-unit eigenvalue, hence eigenvalue \(1\) or \(-1\), contrary to hyperbolicity. Absolute smooth convergence justifies the sums. Equation (ND.12), or (ND.14), now proves (ND.18). End of proof.

For an even finite Fourier polynomial, there is a useful exact identity behind this criterion. If \(a_{\mathcal O}\) denotes the cosine sum of h on a signed orbit, then \[
 \boxed{\sum_{\mathcal O}q_A(v_{\mathcal O})a_{\mathcal O}(h)
                         =-2\pi t\mathcal E_A(M).}          \tag{ND.21}
\] It follows from (ND.13)--(ND.14) using \(\mathcal L_A b_v(0)=-8\pi^2q_A(v)\) and \(\mathcal E_A(-8\pi^2M)=64\pi^4\mathcal E_A(M)\). Thus this test excludes a \emph{nonzero coboundary bracket}, not just an identically zero bracket.

A coarse finite-bandwidth constant is explicit. Define \[
 m_A=\gcd(|c|,|d-a|,|b|),\quad C_A=|c|+|d-a|+|b|,\quad
 D_A=1+\max\{|a|+|b|,|c|+|d|\}.
\] For the one-sign class supported in \(\|v\|_\infty\leq R\), \[
 \boxed{\displaystyle
 \Gamma_A(g)\geq
 \frac{\pi^2t^2m_A^4}{2C_A^2D_A^6}\,R^{-6}\|g\|_2^4.}      \tag{ND.22}
\] Indeed, all generated frequencies lie in the radius-\(D_AR\) box. Choosing representatives from those generated frequencies gives \(\sum_{\mathcal O}q_A(v_{\mathcal O})^2\leq4C_A^2(D_AR)^6\). Cauchy--Schwarz in (ND.21), (ND.1), (ND.20), and \(|q_A(v)|\geq m_A\), \(\|g\|_2^2=2\sum c_v^2\), prove (ND.22). For the Pell matrix this specializes to \[
                \Gamma_T(g)\geq\frac{\pi^2}{32768R^6}\|g\|_2^4.
 \tag{ND.23}
\] No sharpness is claimed. There is no conflict with superpolynomially small positive variances of general Fourier approximations to smooth flat symbols: those approximations need not have one-sign coefficients.

\subsection{An exact finite quantum variance consequence}\label{nondegeneracy:an-exact-finite-quantum-variance-consequence}

The elementary second-moment bridge already occurs in earlier project notes. We reprove it here with an explicit integer sufficient bound so its application to Theorem 1 is auditable without a Gaussian-limit theorem. Let K be any odd positive integer greater than one, write \[
 N+F\sqrt2=(1+\sqrt2)^K,
 \qquad T=\begin{pmatrix}3&2\\4&3\end{pmatrix},
\] and use the physical Weyl operators \[
 Q_N(a,b)=(-1)^{ab}W_N(a,b),\qquad
 Q_N(v)Q_N(w)=e^{\pi i\det(v,w)/N}Q_N(v+w).
\] Here \(W_N\) uses modular division by two and is an orthonormal operator basis in normalized full-space Hilbert--Schmidt norm. Take \(U_N\) with the exact covariance \(U_NQ_N(v)U_N^*=Q_N(Tv)\). Such a unitary is explicitly furnished by the cat kernel in the earlier project; its phase is immaterial here. Primality of K is not needed in this section.

Put \(H=(K-1)/2\) and let \(r_1,r_2\) be the absolute row sums of \(T^H\). For a real mean-zero polynomial h supported in \(\|v\|_\infty\leq M\), assume \[
                    N^2>M^2(4+4r_1+2r_2).                  \tag{ND.24}
\] Define \[
 X_{K,N}(h)=\frac1{\sqrt K}\sum_{j=0}^{K-1}
                    U_N^j\operatorname{Op}_N(h)U_N^{-j}.
\] Then the following equality is \textbf{exact at finite N} on the full space: \[
             \boxed{\frac1N\operatorname{Tr}X_{K,N}(h)^2
                                  =\sigma_T^2(h).}          \tag{ND.25}
\] No parity assumption is made. The equality is not an assertion about individual eigenphase spacings of a kicked propagator.

\textbf{Proof.} The Pell norm gives \(N^2-2F^2=-1\), whence \[
 T^K=\begin{pmatrix}2N^2+1&2NF\\4NF&2N^2+1\end{pmatrix}
                           \equiv I\pmod{2N}.               \tag{ND.26}
\] Thus \(Q_N(T^Kv)=Q_N(v)\), including its physical sign, and adjoint covariance has period K. Differences of times can be reduced to \(-H\leq s\leq H\).

Let \(u,v\) lie in the frequency box and suppose \(u+T^sv\equiv0\pmod N\). Use the indefinite integral norm \(q(a,b)=b^2-2a^2\), invariant under T and nonzero on every nonzero integer vector. Since the row absolute sums of \(T^s\) are bounded by those of \(T^H\), \[
 |q(u+T^sv)|
 \leq |q(u)|+|q(v)|+2|u_2(T^sv)_2-2u_1(T^sv)_1|
 \leq M^2(4+4r_1+2r_2).
\] If \(u+T^sv\) were nonzero, its norm would be a nonzero integer multiple of \(N^2\), contradicting (ND.24). Hence the congruence lifts to the exact relation \(u=-T^sv\) in \(\mathbb Z^2\).

Expand the square in (ND.25). Cyclicity of trace and covariance reduce its normalized value to \[
 \sum_{s=0}^{K-1}\frac1N\operatorname{Tr}
        \bigl(\operatorname{Op}_N(h)\operatorname{Op}_N(\mathcal C_T^s h)\bigr).
\] Weyl traces vanish except at the zero modular frequency. Every surviving relation lifts as above; its two Weyl factors are adjoints and have product I, with no residual phase. There is no double counting of an integer relation since a nonzero integer vector has no nontrivial T-period. Conversely, any integer orbit relation between supported frequencies reduces modulo K and lifts back to the same centered relation. Thus the sum is exactly the oriented-orbit expression (ND.1). End of proof.

For \(g\in\mathcal P_8\), its symbol \(h_g=\{g,\mathcal C_Tg\}/(4\pi)\) is supported in the radius-64 box. At \(K=13\), \(N=47321\), the two sides of the sufficient inequality (ND.24) are the exact integers \[
 N^2=2239277041>935895040=64^2(4+4r_1+2r_2).
\] Therefore (ND.25) and Theorem 1 give a strictly positive quantum averaged-symbol variance for \textbf{every nonzero} such \(g\) at that module, not just the sampled coefficients. The same conclusion holds at every odd K satisfying (ND.24); for each fixed M the condition holds eventually because its right side grows like \(\lambda^{K/2}\) and \(N^2\) grows like \(\lambda^K\).

The actual finite quantum commutator differs from that symbol by a controlled remainder. Put \(G_N=\operatorname{Op}_N(g)\), \(D_N=U_NG_NU_N^*\) and \(R_N=(iN/2)[G_N,D_N]\). The sine remainder gives \[
 \|R_N-\operatorname{Op}_N(h_g)\|_{\rm op}
     \leq\frac{\pi^3}{6N^2}M_3(g)M_3(\mathcal C_Tg),
 \qquad M_3(a)=\sum_v|\widehat a(v)|\,|v|_2^3.
\]
Here \(M_3\) is local to this section (Euclidean frequency weight); it is
not the \(\mathcal M_r\) of the locality section. Consequently, with normalized Hilbert--Schmidt norm, \[
 \left|\left\|\frac1{\sqrt K}\sum_{j=0}^{K-1}U_N^jR_NU_N^{-j}\right\|_{2,N}
                      -\sqrt{\Gamma_T(g)}\right|
 \leq\frac{\pi^3\sqrt K}{6N^2}M_3(g)M_3(\mathcal C_Tg).      \tag{ND.27}
\] This proves nonvanishing of the averaged commutator for each fixed nonzero radius-eight symbol for all sufficiently large admissible N. It does not assert exact equality between the commutator variance and the classical variance at finite N.

Twelve exact Gaussian-integer coefficient checks compare (ND.25) with a direct modular Weyl-coefficient squared norm, including dimensions 47321 and 1607521. No dense matrices of those sizes are constructed. A small-module negative control, \(K=3,N=7\), violates the non-aliasing condition and gives different finite quantum and classical values; this mismatch is retained rather than removing hypothesis (ND.24).
